\chapter{Decoration: One Map, Many Semirings}
\label{sec:weighted}
% ===================================================================

An earlier presentation of this material carried two separate pieces of
data --- a \emph{weight map} valued in $\CC$ and a \emph{cost map} valued
in an ordered monoid --- and called the first a Lagrangian and the second
a Hamiltonian.  That presentation was doing more work than it needed to
and claiming more than it had earned.  The two maps have the same domain,
the same indexing, and the same shape; they differ only in where they take
their values.  This chapter replaces them with a single construction,
parametric in a semiring, and is careful about which classical structures
the construction actually delivers.

\section{Decorations}

\begin{definition}[Semiring]
A \emph{semiring} $\Semi = (R, \oplus, \otimes, 0_R, 1_R)$ is a set with
two monoid structures, $\otimes$ distributing over a commutative $\oplus$,
with $0_R$ absorbing for $\otimes$.  We do not require additive inverses;
that omission is the point of the generality.
\end{definition}

\begin{definition}[Decoration]
\label{def:decoration}
A \emph{decoration} of a GSLT $\GSLT$ in a semiring $\Semi$ is a family
\[
  \dec_r : \HML(\ctx) \rightharpoonup R, \qquad r \in \rules,
\]
whose domain consists of formulae that \emph{refine the left-hand side of
$r$}: formulae $\phi$ such that $u \sat \phi$ implies $u$ matches the
pattern $l_r$.  A \emph{decorated GSLT} is a pair $(\GSLT, \dec)$.
\end{definition}

The refinement relation orders the domain of $\dec_r$: write
$\phi \leq \psi$ when $\psi$ entails $\phi$.  To evaluate a decoration at
a step one must select a formula, and the selection has to be well
defined.

\begin{definition}[Refinement-complete decoration]
\label{def:refcomplete}
A decoration is \emph{refinement-complete} if for every rule $r$ and every
term $u$ matching $l_r$, the set of formulae in $\mathrm{dom}(\dec_r)$
satisfied by $u$ has a greatest element $\phi_r(u)$ --- a most refined
witness.  We then write $\dec_r(u)$ for $\dec_r(\phi_r(u))$.
\end{definition}

This condition was previously left implicit, and it is not automatic: a
domain closed under conjunction with finite refinement classes has it; an
arbitrary subset of $\HML(\ctx)$ does not.  Everything below assumes
refinement-completeness, and the assumption is a real constraint on how
one is permitted to specify a decoration.

\begin{definition}[Path weight]
\label{def:pathweight}
Let $\gamma = r_1 \cdot r_2 \cdots r_n$ be a rewrite path from $P$ to $Q$,
with $u_i$ the redex at step $i$.  The \emph{path weight} is the ordered
product
\[
  W(\gamma) \;=\; \dec_{r_1}(u_1) \otimes \cdots \otimes \dec_{r_n}(u_n)
  \;\in\; R,
\]
and the \emph{transition weight} from $P$ to $Q$ is the sum over paths
\[
  \langle Q \mid P \rangle \;=\; \bigoplus_{\gamma : P \to Q} W(\gamma),
\]
whenever that sum exists.
\end{definition}

The product is ordered because $\otimes$ need not be commutative, and is
not in the cases of interest: the order of the factors is the order of the
steps.  This is the temporal monoid of~\cite[\S4]{CE} reappearing inside
the decoration --- reduction is sequential even when interaction is
parallel.

\section{Four instances}

The content of the parametrisation is that familiar dynamical structures
are recovered by a choice of semiring.  What follows is the reading
of~\cite[\S13]{CE}, transposed to the present setting.

\begin{center}
\renewcommand{\arraystretch}{1.35}
\begin{tabular}{>{\itshape}p{4.3cm} p{7.4cm}}
\toprule
\normalfont\textbf{Semiring $\Semi$} & \textbf{What $(\GSLT,\dec)$ becomes} \\
\midrule
Boolean $(\{0,1\},\vee,\wedge)$
  & the bare nondeterministic transition system; $\langle Q\mid P\rangle$
    is reachability \\
Tropical $(\Real\cup\{\infty\},\min,+)$
  & cost; $W(\gamma)$ is the total charge of a path and
    $\langle Q\mid P\rangle$ is the cheapest way to get there \\
Non-negative reals $(\Real_{\geq 0},+,\times)$
  & stochastic dynamics; $W$ is a path probability, and the Gillespie
    algorithm samples it \\
Complex numbers $(\CC,+,\times)$
  & amplitudes, and with them interference --- see
    Chapter~\ref{sec:quantum}, and the obstruction recorded there \\
\bottomrule
\end{tabular}
\end{center}

\begin{remark}[What the two-map presentation was seeing]
The weight map was the complex instance; the cost map was the tropical
instance.  They were never two kinds of structure.  Nothing is lost by
the merge, and something is gained: the Boolean and stochastic instances,
for which the two-map presentation had no room, are now the same
construction specialised twice more.
\end{remark}

\section{Where the Lagrangian actually is}

The earlier text set $S[\gamma] = \sum_i \log \dec_{r_i}(u_i)$, observed
that $W(\gamma) = e^{S[\gamma]}$, and called $S$ an action.  That is a
change of variable, not a variational principle: it converts a product
into a sum and asserts nothing further.  A Lagrangian earns its name by
making the realised trajectory the stationary point of an action over a
space of competing trajectories, and no such principle was stated, let
alone proved.

There is nevertheless a genuine least-action statement available, and it
is the tropical instance.

\begin{proposition}[Least action is the tropical path integral]
\label{prop:leastaction}
In the tropical semiring, Definition~\ref{def:pathweight} reads
\[
  W(\gamma) = \sum_{i=1}^{n} \dec_{r_i}(u_i),
  \qquad
  \langle Q \mid P \rangle = \min_{\gamma : P \to Q} W(\gamma).
\]
The path weight is an additive action functional, and the transition
weight is its \emph{minimisation over all paths}.  The principle of least
action is therefore not an approximation to the sum over paths in this
setting; it \emph{is} the sum over paths, computed in the semiring whose
addition is $\min$.
\end{proposition}

\begin{proof}
Immediate from $\oplus = \min$ and $\otimes = +$.
\end{proof}

\begin{remark}[And the saddle point]
The relation between Proposition~\ref{prop:leastaction} and the complex
instance is the relation between tropicalisation and the $\hbar \to 0$
limit: the saddle-point approximation of a complex sum over paths is a
tropical minimisation, which is why the classical limit of a path
integral is a least-action principle.  Here the two are not limits of
each other but instances of one construction at different semirings, so
the passage between them is a question about semiring homomorphisms
rather than about asymptotics.  We do not claim to have answered it.  We
claim only that this is where the question lives.
\end{remark}

The honest summary is this.  The decoration supplies a charge per step.
The tropical instance supplies a least-action principle.  The word
``Lagrangian'' is warranted for the tropical case and premature for the
others.

\section{Decorations move}

A decoration as defined above is static: $\dec_r$ is fixed once and for
all.  That is too weak for what the later parts of this book require,
where the rate at which a rule fires depends on what has already happened.

\begin{definition}[Dynamic decoration]
\label{def:dyndec}
A \emph{dynamic decoration} equips each state with a table
$T : \mathrm{dom}(\dec) \to R$ keyed by the formulae refining the
left-hand sides of the rules, and equips each rewrite with an update
$T \mapsto T'$ on the table carried by its successor.  A static decoration
is the special case in which every update is the identity.
\end{definition}

Because the successor state carries a table of the same kind, the updates
compose, and the weights are genuinely dynamical rather than
parameters~\cite{StochSim}.  One consequence is worth recording early:
the adequacy of a logic generated over such a structure is a function of
the rewrites.  Interleaving rewrites yield a logic that sees only
interleaving; true-concurrency rewrites yield a logic that sees true
concurrency.  The decoration does not fix this; the rewrites do.

\section{The cost instance, and affordability}

The tropical instance deserves separate treatment, because it alone
supports a notion of \emph{sufficiency}: a step can fail for want of
funds, which is not something a probability or an amplitude can do.

\begin{definition}[Resource algebra]
A \emph{resource algebra} is a commutative ordered monoid
$(A, +, \mathbf{0}, \leq)$.  A \emph{vectorial account} over resource
types $\mathbf{I} = \{1,\ldots,k\}$ is an element
$\vect{A} = (A_1,\ldots,A_k)$ of $A^k$, each $A_i$ tracking the available
amount of one type.
\end{definition}

For the rho calculus the natural types include a budget for fresh names
(the $@$ constructor), a budget for synchronisations, and a budget for
persistent-receive firings.  Heterogeneous resources are tracked
separately because they are not interconvertible for free; what it costs
to convert one into another is the subject of the next chapter.

\begin{definition}[Account-aware rewrite]
\label{def:affordable}
A step is \emph{affordable} if the account covers the charge:
$\langle P, \vect{A}\rangle \rewrite_r \langle P', \vect{A}'\rangle$
provided that $P \rewrite_r P'$ via $r$ with witness $\phi_r(u)$, that
$\dec_r(u) \leq \vect{A}$ componentwise, and that
$\vect{A}' = \vect{A} - \dec_r(u)$.
\end{definition}

Definition~\ref{def:affordable} is stated here in the form the earlier
presentation used, and in that form it harbours a defect: it says nothing
about \emph{who} may draw on $\vect{A}$.  A single account adjacent to an
associative--commutative soup of processes is adjacent to every process in
it, and adjacency taken as the right to spend is ambient
authority~\cite{finrho}.  The next chapter states the defect precisely and
repairs it, and the repair turns out to carry the most interesting physics
in this part of the book.

% ===================================================================
